\section{A Unified Perspective}
\label{sec:unified_perspective}

The three contributions of this dissertation, while addressing different technical problems, share a common underlying principle: structured abstractions determine how effectively an agent can manage the mismatch between its internal models and the environment. This section develops the connections between the contributions and articulates the unified perspective they support.

\paragraph{Abstraction as a mediator of novelty.}
The conceptual and formal framework developed in \cref{ch:novelty_open_world,ch:foundations} characterized novelty as a mismatch that can affect any component of the agent--environment interaction: objects, attributes, dynamics, or actions. Left unmanaged, such mismatch degrades performance globally---the agent's entire policy may become ineffective, and recovery requires extensive relearning. The contributions of this dissertation demonstrate that structured abstractions can transform this global problem into a collection of localized ones.

Action abstractions achieve this at the level of task structure. By decomposing behavior into a sequence of operators with explicit preconditions and effects, each associated with a separate executor, the agent creates natural interfaces at which mismatch can be detected and isolated. When novelty causes a specific executor to fail, the agent knows \emph{which} component is affected and \emph{what} the replacement must achieve, reducing adaptation from a full policy search to a targeted learning problem. Interaction abstractions achieve this at the level of individual behaviors. By representing policies in force space rather than robot-centric space, the agent factors out the embodiment-specific details that vary across platforms and objects, ensuring that the learned behavior captures only the task-relevant physical invariants. When a new object or robot is encountered, the invariant representation means that transfer requires only a mapping between force commands and the specific embodiment, rather than relearning the task from scratch.

\paragraph{Complementary roles of assimilation and accommodation.}
The distinction developed in \cref{ch:novelty_open_world} provides a common interpretation of the adaptation and generalization contributions. Accommodation is required when existing competence is insufficient: the agent must diagnose mismatch and revise or extend the affected knowledge or behavior. Assimilation handles variation within an existing representation; in the generalization chapters, invariant interaction representations allow an existing policy to remain applicable without structural revision.

The contributions illustrate both responses and their interplay. Action abstractions support localized accommodation (\cref{ch:rapid_learn}): when novelty invalidates an operator or executor, structured decomposition enables targeted recovery while preserving unaffected competence. Interaction abstractions support assimilation through generalization (\cref{ch:flex,ch:pho}): force-space representations transfer across the evaluated objects, physical conditions, and embodiments without policy retraining. Neither response alone is sufficient for open-world operation---accommodation without generalization requires repeated revision for every variation, while assimilation within a fixed representation provides no recourse when that representation becomes insufficient. The unified perspective is that a robust open-world agent requires both, organized through abstractions that determine which variations are handled within invariant representations and which require active structural or behavioral revision.

\paragraph{Toward integrated architectures.}
As discussed in \cref{sec:perspectives_integration}, the natural integration of these contributions would produce an agent architecture in which the planning layer provides task decomposition and failure localization (action abstractions), the execution layer provides generalizable force-based skills (interaction abstractions), and the adaptation layer provides targeted learning when execution fails. The benchmarks developed in \cref{ch:benchmarks} would provide the evaluation infrastructure for such an integrated system. While this full integration was not realized within this dissertation, the formal framework of \cref{ch:foundations} and the empirical results of \cref{ch:rapid_learn,ch:flex,ch:pho} establish the components and demonstrate their individual effectiveness. The integration represents the most immediate direction for future work.
